\section{Introduction}

\paragraph{Problème.} Retrouver la règle applicable dans les codes tunisiens demande de savoir dans quel code chercher,
sous quel terme juridique la règle est formulée, et dans quel article précis elle figure. Les codes sont publiés en
français ; les questions des usagers arrivent en arabe, en français ou en anglais, souvent dans un vocabulaire courant
(« frapper quelqu'un ») éloigné de celui du texte (« coups et blessures volontaires »). Un système de génération augmentée
par récupération (\term{RAG}) répond à ce besoin en récupérant des articles puis en rédigeant une réponse à partir d'eux.
Son point faible est double : il peut ne pas trouver le bon article, et il peut affirmer ce que les articles ne disent pas.

\paragraph{Objectifs.} Nous poursuivons trois objectifs, dans cet ordre :
\begin{enumerate}
  \item mesurer et corriger la récupération du système existant (v1), pour obtenir une référence honnête (v1.1) ;
  \item construire un agent (v2) qui décide lui-même quand chercher davantage, et dont chaque affirmation cite un article
        vérifiable ;
  \item évaluer v1, v1.1 et v2 sur les mêmes critères, en séparant strictement les données de mise au point et
        les données de test.
\end{enumerate}

\paragraph{Périmètre.} Le corpus comprend \NCorpusPdfs{} codes tunisiens au format PDF (\NCorpusPages{} pages,
\NCorpusChunks{} \term{chunks} dans l'index). Le jeu de questions de récupération compte \NQTotal{} questions
(\NQLangFr{} en français, \NQLangAr{} en arabe, \NQLangEn{} en anglais), groupées en \NConceptsTotal{} concepts juridiques.
Nous évaluons la récupération (l'article de référence est-il trouvé ?) et, pour l'agent, la génération (la réponse
est-elle exacte et appuyée par les articles cités ?). L'évaluation au niveau de la réponse sur un jeu de cent questions
n'est pas encore réalisée ; la section~\ref{sec:finale} en donne le protocole.

\paragraph{Plan.} La section~\ref{sec:corpus} décrit le corpus et la vérification des \term{chunks} contre les PDF ;
la section~\ref{sec:v11} le passage de v1 à v1.1 ; la section~\ref{sec:methodo} la méthodologie d'évaluation ;
la section~\ref{sec:archi} l'architecture de l'agent ; les sections~\ref{sec:recup} à~\ref{sec:stabilite} les expériences ;
la section~\ref{sec:menaces} les menaces à la validité.
